\usepackage{array}

% Table:
\usepackage{booktabs} % 用于绘制更美观的表格线

% Math:
\usepackage{sansmath}

% 水印:
% \usepackage{eso-pic}
% \AddToShipoutPictureBG{ % 这会将内容添加到页面的背景
%     \put(0,0){ % put 命令用于放置内容，(0,0) 是左下角
%         \parbox{\paperwidth}{\centering % parbox 设置一个与页面同宽的盒子，并居中对齐
%             \color[rgb]{0.9,0.9,0.9} % 设置文本颜色为浅灰色
%             \fontsize{5cm}{5cm}\selectfont % 设置字体大小
%             By Shichien % 这是水印文本
%             % \rotatebox{45}{\textsf{\textbf{By Shichien}}} % 如果需要旋转水印，可以使用这个
%         }
%     }
% }

% Drawing
\usepackage{pgfplots}
\usepackage{tikz-cd}
\usepackage{tikz-qtree}
\usepackage{circuitikz}
\usepackage{tikz-3dplot}
\usepackage[svgnames]{xcolor}
\usetikzlibrary{shapes.geometric, arrows.meta, positioning, calc, angles, quotes, intersections, tikzmark}

\pgfplotsset{compat=1.18}
\newcommand{\ccr}[1]{\makecell{{\color{#1}\rule{1cm}{1cm}}}}
\usepackage{algorithm,algorithmic}

% 选择题选项环境
\usepackage{tasks}
\settasks{label = \Alph*.
    %   label-offset={0.4em},
    %     label-align=left,
    %     column-sep={2pt},
    %     item-indent={1pt},before-skip={-0.7em},after-skip={-0.7em}
}

% TColorbox 
\usepackage{tcolorbox}
\usepackage{fontspec}
\usepackage{fontawesome5}

\definecolor{warningRed}{RGB}{217, 83, 79}
\definecolor{warningRedBack}{RGB}{242, 222, 222}

\newtcolorbox{warning}{
    enhanced,
    shadow={0.8mm}{-0.8mm}{0mm}{black!20!white},
    breakable, % 允许内容跨页
    colback=warningRedBack, % 背景颜色
    colframe=warningRed,    % 边框颜色
    boxrule=0.5pt,          % 边框线宽
    leftrule=2pt,           % 左侧竖线的宽度 (比边框粗，更明显)
    top=2pt,
    bottom=1.8pt,
    right=5em,
    fontupper=\fontsize{9.5pt}{13pt}\selectfont\kaishu,
    fonttitle=\small\bfseries,
    title={
            \faExclamationTriangle[regular]
            \hspace{1.5pt}
            \textbf{Warning}
        },
    coltitle=white,
    attach boxed title to top right={xshift=0pt, yshift=-18.3pt},
    boxed title style={
            colback=warningRed,
            boxrule=0.3pt,
        }
}

\definecolor{tipBlue}{RGB}{91, 192, 222}
\definecolor{tipBlueBack}{RGB}{217, 237, 247}

\newtcolorbox{tips}{
    enhanced,
    shadow={0.8mm}{-0.8mm}{0mm}{black!20!white},
    breakable,
    colback=tipBlueBack,    % 背景颜色
    colframe=tipBlue,       % 边框颜色
    boxrule=0.5pt,          % 边框线宽
    leftrule=2pt,           % 左侧竖线的宽度
    top=2pt,
    bottom=1.8pt,
    right=5em,
    fontupper=\fontsize{9.5pt}{13pt}\selectfont\kaishu, % 内容字体: {大小}{衬线}
    fonttitle=\small\bfseries,
    title={
            \faLightbulb[regular]
            \hspace{1.5pt}
            \textbf{Tips}
        },
    coltitle=white,
    attach boxed title to top right={xshift=0pt, yshift=-18.3pt},
    % attach boxed title to top left={xshift=10pt, yshift=-10pt},
    boxed title style={
            colback=tipBlue,
            boxrule=0.3pt,
        }
}

\newcommand{\h}{\hspace}
\renewcommand{\v}{\vspace}

% 页脚修改，加入 Shichien 的 Github 图标和链接
\usepackage{fancyhdr}
\pagestyle{fancy}
\fancyfoot[L]{%
    \href{https://github.com/Shichien}{\faGithub\ \textit{Shichien}}
}
\fancyfoot[C]{\thepage} % 页脚中间（C）显示页码
\renewcommand{\headrulewidth}{1pt} % 保留页眉分割线
\renewcommand{\footrulewidth}{0.4pt} % 添加页脚分割线（可选）

% 代码块设置
\definecolor{lstcolorComment}{rgb}{0,0.6,0}
\definecolor{lstcolorNumber}{HTML}{E6E6E6}
\definecolor{lstcolorString}{rgb}{1,0.5,0}

\usepackage{listings}
\lstset{
    backgroundcolor=\color{white},           % 背景色
    basicstyle=\footnotesize\ttfamily,       % 代码字体和大小 (增加 ttfamily 保证等宽)
    breakatwhitespace=false,                 % 是否只在空白处断行
    breaklines=true,                         % 自动断行
    keepspaces=true,                         % 保留代码中的空格，对齐代码
    tabsize=2,                               % Tab 长度为 2 个空格
    frame=single,                            % 单线边框
    rulecolor=\color{black},                 % 边框颜色
    captionpos=bl,                           % 标题位置 (bottom-left)
    numbers=left,                            % 行号在左侧
    numbersep=5pt,                           % 行号与代码的距离
    numberstyle=\tiny\color{lstcolorNumber}, % 行号样式
    stepnumber=1,                            % 行号步长
    commentstyle=\color{lstcolorComment},    % 注释样式
    keywordstyle=\color{blue},               % 关键字样式
    stringstyle=\color{lstcolorString},      % 字符串样式
    showspaces=false,                        % 不显示空格
    showstringspaces=false,                  % 不显示字符串中的空格
    showtabs=false,                          % 不显示 Tab
    escapeinside={\%*}{*},                     % 允许在代码中插入 LaTeX 命令
    extendedchars=true                       % 支持非 ASCII 字符 (对旧编码有效)
}

% % Colour Palette ——————————————————————————————————————
% \definecolor{merah}{HTML}{F4564E}
% \definecolor{merahtua}{HTML}{89313E}
% \definecolor{biru}{HTML}{60BBE5}
% \definecolor{birutua}{HTML}{412F66}
% \definecolor{hijau}{HTML}{59CC78}
% \definecolor{hijautua}{HTML}{366D5B}
% \definecolor{kuning}{HTML}{FFD56B}
% \definecolor{jingga}{HTML}{FBA15F}
% \definecolor{ungu}{HTML}{8C5FBF}
% \definecolor{lavender}{HTML}{CBA5E8}
% \definecolor{merjamb}{HTML}{FFB6E0}
% \definecolor{mygray}{HTML}{E6E6E6}
% \definecolor{mygreen}{rgb}{0,0.6,0}
% \definecolor{mymauve}{rgb}{0.58,0,0.82}

% % Theorems ————————————————————————————————————————————
% \usepackage{tcolorbox}
% \usepackage{changepage}
% \tcbuselibrary{skins,breakable,theorems}

% \makeatletter
% % 定义一个标题样式，格式为： "环境名 环境编号 (可选标题)"
% % #1=环境名, #2=可选标题, #3=编号
% \def\tcb@theo@widetitle#1#2#3{%
%     \hbox to \textwidth{\bfseries\large#1\ \normalsize#3\if\relax\detokenize{#2}\relax\else\space(#2)\fi\hfil}%
% }
% \tcbset{
%     theorem style/mywide/.style={
%             let\tcb@theo@title=\tcb@theo@widetitle
%         }
% }
% \makeatother


% % ===================================================================
% % 2. 环境定义
% % ===================================================================

% % -- 创建一个共享的计数器 --
% \newcounter{hitung}

% % % -- 证明 (Proof) --
% % \newtcolorbox{proof}[1][证明]{
% %     enhanced, breakable, skin=enhancedmiddle,
% %     parbox=false, boxrule=0mm,
% %     % 样式
% %     colframe=black!25,
% %     colback=black!5,
% %     leftrule=3pt,
% %     % 间距
% %     boxsep=0mm, left=3mm, right=3mm, top=2mm, bottom=2mm,
% %     before skip=8pt, after skip=8pt,
% %     % 标题
% %     fonttitle=\bfseries, coltitle=black,
% %     title=#1,
% %     % 结束时自动添加证毕符号
% %     finish={\null\hfill$\blacksquare$},
% % }

% % -- 定义 (Definition) & 解答 (Answer) --
% \newtcbtheorem[use counter=hitung, number within=section]{dfn}{定义}{
%     theorem style=mywide, breakable, enhanced,
%     arc=3.5mm, outer arc=3.5mm,
%     boxrule=0pt, toprule=1pt, bottomrule=1pt,
%     left=2mm, right=2mm, top=2mm, bottom=2mm,
%     toptitle=1mm, bottomtitle=1mm,
%     before skip=8pt, after skip=8pt,
%     % 颜色与字体 (蓝色系)
%     colframe=blue!75!black,
%     colback=blue!10,
%     coltitle=white,
%     title style={colback=blue!75!black},
%     fonttitle=\bfseries,
%     coltext=blue!50!black,
%     % 阴影
%     shadow={1mm}{-1mm}{0mm}{gray!50},
% }{dfn}

% \newtcbtheorem[use counter=hitung, number within=section]{ans}{解答}{
%     % 样式与“定义”环境完全相同
%     theorem style=mywide, breakable, enhanced,
%     arc=3.5mm, outer arc=3.5mm,
%     boxrule=0pt, toprule=1pt, bottomrule=1pt,
%     left=2mm, right=2mm, top=2mm, bottom=2mm,
%     toptitle=1mm, bottomtitle=1mm,
%     before skip=8pt, after skip=8pt,
%     colframe=blue!75!black, colback=blue!10, coltitle=white,
%     title style={colback=blue!75!black}, fonttitle=\bfseries,
%     coltext=blue!50!black, shadow={1mm}{-1mm}{0mm}{gray!50},
% }{ans}

% % -- 公理 (Axiom) & 命题 (Proposition) --
% \newtcbtheorem[use counter=hitung, number within=section]{axm}{公理}{
%     theorem style=mywide, breakable, enhanced,
%     arc=3.5mm, outer arc=3.5mm,
%     boxrule=0pt, toprule=1pt, bottomrule=1pt,
%     left=2mm, right=2mm, top=2mm, bottom=2mm,
%     toptitle=1mm, bottomtitle=1mm,
%     before skip=8pt, after skip=8pt,
%     % 颜色与字体 (绿色系)
%     colframe=green!60!black,
%     colback=green!10,
%     coltitle=white,
%     title style={colback=green!60!black},
%     fonttitle=\bfseries,
%     coltext=green!40!black,
%     shadow={1mm}{-1mm}{0mm}{gray!50},
% }{axm}

% \newtcbtheorem[use counter=hitung, number within=section]{prp}{命题}{
%     % 样式与“公理”环境完全相同
%     theorem style=mywide,
%     breakable,
%     enhanced,
%     arc=3.5mm,
%     outer arc=3.5mm,
%     boxrule=0pt,
%     toprule=1pt,
%     bottomrule=1pt,
%     left=2mm,
%     right=2mm,
%     top=2mm,
%     bottom=2mm,
%     toptitle=1mm,
%     bottomtitle=1mm,
%     before skip=8pt,
%     after skip=8pt,
%     colframe=green!60!black,
%     colback=green!10,
%     coltitle=white,
%     title style={colback=green!60!black},
%     fonttitle=\bfseries,
%     coltext=green!40!black,
%     shadow={1mm}{-1mm}{0mm}{gray!50},
% }{prp}

% % -- 定理 (Theorem) --
% \newtcbtheorem[use counter=hitung, number within=section]{thm}{定理}{
%     theorem style=mywide,
%     breakable,
%     enhanced,
%     arc=3.5mm,
%     outer arc=3.5mm,
%     boxrule=0pt,
%     toprule=1pt,
%     bottomrule=1pt,
%     left=2mm,
%     right=2mm,
%     top=2mm,
%     bottom=2mm,
%     toptitle=1mm,
%     bottomtitle=1mm,
%     before skip=8pt,
%     after skip=8pt,
%     % 颜色与字体 (红色系)
%     colframe=red!75!black,
%     colback=red!10,
%     coltitle=white,
%     title style={colback=red!75!black},
%     fonttitle=\bfseries,
%     coltext=red!60!black,
%     shadow={1mm}{-1mm}{0mm}{gray!50},
% }{thm}

% % -- 示例 (Example) --
% \newtcolorbox[use counter=hitung, number within=section]{cth}[1][]{
%     breakable,
%     enhanced,
%     % 颜色 (橙色系)
%     colback=orange!15,
%     colframe=orange!80!black,
%     coltitle=white,
%     colbacktitle=orange!80!black,
%     coltext=black!75,
%     fonttitle=\bfseries,
%     before skip=8pt,
%     after skip=8pt,
%     attach boxed title to top left={yshift=-2mm},
%     title={示例 \thetcbcounter\ifdefempty{#1}{}{ (#1)}},
% }

% % -- 注释 (Note) --
% \newtcolorbox{ctt}[1][]{
%     enhanced,
%     breakable,
%     % 关键样式：左侧竖线 (黄色系)
%     borderline west={3pt}{0pt}{yellow!80!black},
%     left=5mm,
%     colback=yellow!20,
%     colframe=yellow!20,
%     coltext=yellow!70!black,
%     coltitle=yellow!60!black,
%     fonttitle=\bfseries\scshape,
%     before skip=8pt,
%     after skip=8pt,
%     title=注释\ifdefempty{#1}{}{ (#1)},
% }

% % -- 备注 (Remark) --
% \newtcolorbox{rmr}[1][]{
%     breakable,
%     enhanced,
%     % 关键样式：底部横线 (灰色系)
%     arc=0mm,
%     boxrule=0pt,
%     bottomrule=2pt,
%     colback=black!8,
%     colframe=black!60,
%     coltext=black!80,
%     fonttitle=\bfseries,
%     top=3mm,
%     bottom=3mm,
%     before skip=8pt,
%     after skip=8pt,
%     title=备注\ifdefempty{#1}{}{ (#1)},
% }